% GENERATED FILE. Edit canonical lessons, then run npm run generate:books.
% canonical-source-hash: fnv1a64:0efe76dbb4a4a7ce
% canonical-lessons: TA-C472-catu, TA-C472-vekuli, TA-C472-emali, TA-C472-kancan, TA-C472-utari

\chapter{Kinds of People}
\label{ch:ta-kinds-of-people}
\hlchaptermodality{472}

\begin{chapteropening}
\textbf{By the end of this chapter:} \emph{I can describe gentle, naive, gullible, stingy and wasteful people.}

Complete the last lesson of chapter 472: I can describe gentle, naive, gullible, stingy and wasteful people.
\end{chapteropening}

\section[\texorpdfstring{\ta{சாது} (cātu)}{சாது (cātu)}]{\ta{சாது} (cātu) — a meek person; a gentle soul}

\label{lesson:TA-C472-catu}



\begin{quote}
Before the new one: say the Tamil for maturity, then the Tamil for negligence.
\end{quote}

\subsection*{You'll want to know: \ta{சாது}}
\textbf{\ta{சாது}} — \emph{cātu} — \textquotedblleft{}a meek person; a gentle soul\textquotedblright{}.

\begin{quote}
\ta{என்} \ta{தாத்தா} \ta{ஒரு} \ta{சாது}, \ta{அவருக்குக்} \ta{கோபமே} \ta{வராது}.
\end{quote}

\begin{quote}
\emph{eṉ tāttā oru cātu, avarukkuk kōpamē varātu.}
\end{quote}

\begin{quote}
My grandfather is a gentle soul; he never gets angry.
\end{quote}

Also a holy man; \ta{சாதுவான} \ta{மாடு}, a docile cow.

\begin{cousinweb}
\textbf{\ta{சாது}} (\emph{cātu}) is borrowed from Sanskrit \emph{sādhu}, \textquotedblleft{}straight, good; a good man\textquotedblright{} (CDIAL 13337).
\end{cousinweb}

\begin{grammarlens}[title={What you've built}]
One more word for describing what you lived and felt.
\end{grammarlens}

\subsection*{Your turn}
\begin{itemize}
\raggedright
  \item \emph{Say it:} \emph{cātu}
  \item \emph{Say it:} \emph{cātu}, once more
  \item \emph{Say it:} \emph{eṉ tāttā oru cātu, avarukkuk kōpamē varātu.} — My grandfather is a gentle soul; he never gets angry.
  \item \emph{Recall:} say the Tamil for maturity, then the Tamil for negligence, then say \emph{cātu} again
\end{itemize}

\begin{tcolorbox}[breakable,colback=teal!4,colframe=teal!35!black,title={Before you move on}]
What does \ta{சாது} mean? (\textquotedblleft{}A meek person; a gentle soul\textquotedblright{}.) Say it once more.
\end{tcolorbox}
\section[\texorpdfstring{\ta{வெகுளி} (vekuḷi)}{வெகுளி (vekuḷi)}]{\ta{வெகுளி} (vekuḷi) — a naive person; an innocent}

\label{lesson:TA-C472-vekuli}



\begin{quote}
Before the new one: say the Tamil for negligence, then the Tamil for a meek person.
\end{quote}

\subsection*{You'll want to know: \ta{வெகுளி}}
\textbf{\ta{வெகுளி}} — \emph{vekuḷi} — \textquotedblleft{}a naive person; an innocent\textquotedblright{}.

\begin{quote}
\ta{அவள்} \ta{ஒரு} \ta{வெகுளி}, \ta{எல்லோரையும்} \ta{நம்புவாள்}.
\end{quote}

\begin{quote}
\emph{avaḷ oru vekuḷi, ellōraiyum nampuvāḷ.}
\end{quote}

\begin{quote}
She is naive; she trusts everyone.
\end{quote}

\ta{வெகுளித்தனம்}, naivety.

\begin{cousinweb}
\textbf{\ta{வெகுளி}} (\emph{vekuḷi}) is native Dravidian: DEDR 5464 gives it the older sense \textquotedblleft{}anger\textquotedblright{}. How it came to mean an artless person is not settled.
\end{cousinweb}

\begin{grammarlens}[title={What you've built}]
One more word for describing what you lived and felt.
\end{grammarlens}

\subsection*{Your turn}
\begin{itemize}
\raggedright
  \item \emph{Say it:} \emph{vekuḷi}
  \item \emph{Say it:} \emph{vekuḷi}, once more
  \item \emph{Say it:} \emph{avaḷ oru vekuḷi, ellōraiyum nampuvāḷ.} — She is naive; she trusts everyone.
  \item \emph{Recall:} say the Tamil for negligence, then the Tamil for a meek person, then say \emph{vekuḷi} again
\end{itemize}

\begin{tcolorbox}[breakable,colback=teal!4,colframe=teal!35!black,title={Before you move on}]
What does \ta{வெகுளி} mean? (\textquotedblleft{}A naive person; an innocent\textquotedblright{}.) Say it once more.
\end{tcolorbox}
\section[\texorpdfstring{\ta{ஏமாளி} (ēmāḷi)}{ஏமாளி (ēmāḷi)}]{\ta{ஏமாளி} (ēmāḷi) — a gullible person; a dupe}

\label{lesson:TA-C472-emali}



\begin{quote}
Before the new one: say the Tamil for a meek person, then the Tamil for a naive person.
\end{quote}

\subsection*{You'll want to know: \ta{ஏமாளி}}
\textbf{\ta{ஏமாளி}} — \emph{ēmāḷi} — \textquotedblleft{}a gullible person; a dupe\textquotedblright{}.

\begin{quote}
\ta{நான்} \ta{ஏமாளி} \ta{இல்லை}, \ta{உன்} \ta{பொய்யை} \ta{நம்ப} \ta{மாட்டேன்}.
\end{quote}

\begin{quote}
\emph{nāṉ ēmāḷi illai, uṉ poyyai nampa māṭṭēṉ.}
\end{quote}

\begin{quote}
I am no fool; I will not believe your lie.
\end{quote}

From the same root as \ta{ஏமாற்று} \textquotedblleft{}to cheat\textquotedblright{}; -\ta{ஆளி} means one who is.

\begin{cousinweb}
\textbf{\ta{ஏமாளி}} (\emph{ēmāḷi}) is native Dravidian (DEDR 898), in the entry of \ta{ஏமாறு} \textquotedblleft{}to be deceived\textquotedblright{}.
\end{cousinweb}

\begin{grammarlens}[title={What you've built}]
One more word for describing what you lived and felt.
\end{grammarlens}

\subsection*{Your turn}
\begin{itemize}
\raggedright
  \item \emph{Say it:} \emph{ēmāḷi}
  \item \emph{Say it:} \emph{ēmāḷi}, once more
  \item \emph{Say it:} \emph{nāṉ ēmāḷi illai, uṉ poyyai nampa māṭṭēṉ.} — I am no fool; I will not believe your lie.
  \item \emph{Recall:} say the Tamil for a meek person, then the Tamil for a naive person, then say \emph{ēmāḷi} again
\end{itemize}

\begin{tcolorbox}[breakable,colback=teal!4,colframe=teal!35!black,title={Before you move on}]
What does \ta{ஏமாளி} mean? (\textquotedblleft{}A gullible person; a dupe\textquotedblright{}.) Say it once more.
\end{tcolorbox}
\section[\texorpdfstring{\ta{கஞ்சன்} (kañcaṉ)}{கஞ்சன் (kañcaṉ)}]{\ta{கஞ்சன்} (kañcaṉ) — a miser}

\label{lesson:TA-C472-kancan}



\begin{quote}
Before the new one: say the Tamil for a naive person, then the Tamil for a gullible person.
\end{quote}

\subsection*{You'll want to know: \ta{கஞ்சன்}}
\textbf{\ta{கஞ்சன்}} — \emph{kañcaṉ} — \textquotedblleft{}a miser\textquotedblright{}.

\begin{quote}
\ta{அந்தப்} \ta{பணக்காரன்} \ta{ஒரு} \ta{கஞ்சன்}, \ta{யாருக்கும்} \ta{உதவ} \ta{மாட்டான்}.
\end{quote}

\begin{quote}
\emph{antap paṇakkāraṉ oru kañcaṉ, yārukkum utava māṭṭāṉ.}
\end{quote}

\begin{quote}
That rich man is a miser; he will help nobody.
\end{quote}

\ta{கஞ்சத்தனம்}, stinginess.

\begin{cousinweb}
The origin of \textbf{\ta{கஞ்சன்}} (\emph{kañcaṉ}) is not settled; it has the -\ta{அன்} ending of a person, but the stem is not in our sources.
\end{cousinweb}

\begin{grammarlens}[title={What you've built}]
One more word for describing what you lived and felt.
\end{grammarlens}

\subsection*{Your turn}
\begin{itemize}
\raggedright
  \item \emph{Say it:} \emph{kañcaṉ}
  \item \emph{Say it:} \emph{kañcaṉ}, once more
  \item \emph{Say it:} \emph{antap paṇakkāraṉ oru kañcaṉ, yārukkum utava māṭṭāṉ.} — That rich man is a miser; he will help nobody.
  \item \emph{Recall:} say the Tamil for a naive person, then the Tamil for a gullible person, then say \emph{kañcaṉ} again
\end{itemize}

\begin{tcolorbox}[breakable,colback=teal!4,colframe=teal!35!black,title={Before you move on}]
What does \ta{கஞ்சன்} mean? (\textquotedblleft{}A miser\textquotedblright{}.) Say it once more.
\end{tcolorbox}
\section[\texorpdfstring{\ta{ஊதாரி} (ūtāri)}{ஊதாரி (ūtāri)}]{\ta{ஊதாரி} (ūtāri) — a spendthrift}

\label{lesson:TA-C472-utari}



\begin{quote}
Before the new one: say the Tamil for a gullible person, then the Tamil for a miser.
\end{quote}

\subsection*{You'll want to know: \ta{ஊதாரி}}
\textbf{\ta{ஊதாரி}} — \emph{ūtāri} — \textquotedblleft{}a spendthrift\textquotedblright{}.

\begin{quote}
\ta{அவன்} \ta{ஒரு} \ta{ஊதாரி}, \ta{சம்பளம்} \ta{ஒரு} \ta{வாரத்தில்} \ta{தீர்ந்துவிடும்}.
\end{quote}

\begin{quote}
\emph{avaṉ oru ūtāri, campaḷam oru vārattil tīrntuviṭum.}
\end{quote}

\begin{quote}
He is a spendthrift; his salary is gone in a week.
\end{quote}

\ta{ஊதாரித்தனம்}, wastefulness.

\begin{cousinweb}
The origin of \textbf{\ta{ஊதாரி}} (\emph{ūtāri}) is not settled; none of our sources lists it.
\end{cousinweb}

\begin{grammarlens}[title={What you've built}]
One more word for describing what you lived and felt.
\end{grammarlens}

\subsection*{Your turn}
\begin{itemize}
\raggedright
  \item \emph{Say it:} \emph{ūtāri}
  \item \emph{Say it:} \emph{ūtāri}, once more
  \item \emph{Say it:} \emph{avaṉ oru ūtāri, campaḷam oru vārattil tīrntuviṭum.} — He is a spendthrift; his salary is gone in a week.
  \item \emph{Recall:} say the Tamil for a gullible person, then the Tamil for a miser, then say \emph{ūtāri} again
\end{itemize}

\begin{tcolorbox}[breakable,colback=teal!4,colframe=teal!35!black,title={Before you move on}]
What does \ta{ஊதாரி} mean? (\textquotedblleft{}A spendthrift\textquotedblright{}.) Say it once more.
\end{tcolorbox}
